\chapter{Fundamentals Ⅱ: Linear Optical Respond}
\label{cha:fundamentals_b}

\section{Toward electromagnetic response in matter}

The previous chapter focused on the ground-state electronic problem.
There, the central question was how electrons arrange themselves under a given external potential.
However, optical properties do not follow from the ground-state alone.
When an electromagnetic field is applied, it drives the electrons away from their static equilibrium and triggers a series of responses.

Therefore, before introducing the dielectric function, we need a language that describes both the applied field and the feedback of matter to that field.
That language is electrodynamics.
The present section does not aim to review a complete lecture of electrodynamics.
Instead, we only extract the ingredients required for the later discussion of the dielectric function, the Kramers--Kronig relations, and the optical properties.
We begin with electromagnetic fields in matter, and then gradually move toward a frequency-dependent response function that can be connected to the electronic structure.

In the present section, we mainly use SI units for the macroscopic electrodynamic equations.
This choice keeps the structure of the Maxwell equations explicit.

\subsection{Maxwell equations in matter}

Let us start with the universal structure of the problem.
When an electromagnetic field enters a material, it drives charge redistribution and induces currents inside the medium, while these induced sources further modify the field itself.
Therefore, the optical response of matter should be described within the framework of electrodynamics~\cite{griffiths2023introduction}.

At this stage, we do not explore how a specific material responds.
Instead, we first identify the general field equations that govern electromagnetic dynamics in matter.
They provide the universal backbone of the problem, while the constitutive relations later encode the material-specific information.

Firstly, we write the macroscopic Maxwell equations for the electric field $\vec{E}(\vec{r},t)$ and magnetic field $\vec{B}(\vec{r},t)$ as:

\begin{equation}
\left\{\begin{aligned}
    \,\nabla\cdot\vec{E}(\vec{r},t) &= \frac{\rho(\vec{r},t)}{\varepsilon_0} 
    &&\quad \text{Gauss's law for electric field}\\
    \,\nabla\cdot\vec{B}(\vec{r},t) &= 0 
    &&\quad \text{Gauss's law for magnetic field}\\
    \,\nabla\times\vec{E}(\vec{r},t) &= -\frac{\partial}{\partial t}\vec{B}(\vec{r},t) 
    &&\quad \text{Faraday's law}\\
    \,\nabla\times\vec{B}(\vec{r},t) &= \mu_0\vec{J}(\vec{r},t)+\mu_0\varepsilon_0\frac{\partial}{\partial t}\vec{E}(\vec{r},t)
    &&\quad \text{Ampère--Maxwell law}
\end{aligned}\right.
\label{maxwell_equations_with_total_sources}
\end{equation}

here, $\rho(\vec{r},t)$ and $\vec{J}(\vec{r},t)$ denote the total charge density and total current density, respectively.

These equations already imply a local conservation law.
More specifically, we take the divergence of the Ampère--Maxwell law,
and use Gauss's law for the electric field,
we obtain the charge continuity equation:

\begin{equation}
    \frac{\partial}{\partial t}\rho(\vec{r},t)+\nabla\cdot\vec{J}(\vec{r},t)=0
\label{macroscopic_charge_continuity_equation}
\end{equation}

% draft
This is the continuity equation for charge at the macroscopic level.
In this sense, the conservation principle encountered in the previous chapter reappears here in electrodynamic language.

Next, in a material it is useful to separate the sources into free and bound parts:
\begin{equation}
\begin{aligned}
    \rho(\vec{r},t)
    &= \rho_\mathrm{free}(\vec{r},t)+\rho_\mathrm{bound}(\vec{r},t) \\
    \vec{J}(\vec{r},t)
    &= \vec{J}_\mathrm{free}(\vec{r},t)+\vec{J}_\mathrm{bound}(\vec{r},t)
\end{aligned}
\label{free_and_bound_source_decomposition}
\end{equation}

The free part describes externally supplied or mobile macroscopic sources, while the bound part describes the response generated internally by the medium.
To encode the bound electric response, we introduce the polarization field $\vec{P}(\vec{r},t)$ through~\cite{jackson1999classical,ashcroft1976solid}

\begin{equation}
    \rho_\mathrm{bound}(\vec{r},t)\coloneqq -\nabla\cdot\vec{P}(\vec{r},t)
\label{bound_charge_from_polarization}
\end{equation}

Accordingly, we define the electric displacement field as:

\begin{equation}
    \vec{D}(\vec{r},t)\coloneqq \varepsilon_0\vec{E}(\vec{r},t)+\vec{P}(\vec{r},t)
\label{electric_displacement_field_definition}
\end{equation}

For the nonmagnetic materials considered in the present work, the electric response is the central object.
Therefore, it is sufficient here to keep the Maxwell equations in the form

\begin{equation}
\left\{
\begin{aligned}
    \nabla\cdot\vec{D}(\vec{r},t) &= \rho_\mathrm{free}(\vec{r},t) \\
    \nabla\cdot\vec{B}(\vec{r},t) &= 0 \\
    \nabla\times\vec{E}(\vec{r},t) &= -\frac{\partial}{\partial t}\vec{B}(\vec{r},t) \\
    \frac{1}{\mu_0}\nabla\times\vec{B}(\vec{r},t) &= \vec{J}_\mathrm{free}(\vec{r},t)+\frac{\partial}{\partial t}\vec{D}(\vec{r},t)
\end{aligned}
\right.
\label{maxwell_equations_in_matter}
\end{equation}

Consequently, the structure of the problem becomes clear.
The Maxwell equations themselves are universal, while the material dependence is carried by the induced response hidden in $\vec{P}(\vec{r},t)$, and therefore in $\vec{D}(\vec{r},t)$.

In many optical problems for bulk solids, one considers the response to an external electromagnetic perturbation in the absence of free charges and free currents inside the medium.
Under this condition, equation~\eqref{maxwell_equations_in_matter} shows that the central issue is no longer the field equations themselves, but rather how the medium builds up polarization under the applied electric field.

Finally, this observation already points to the route that we need.
Once the relation between $\vec{P}(\vec{r},t)$ and $\vec{E}(\vec{r},t)$ is specified, the electromagnetic response of the medium becomes well defined.
This is precisely the point where macroscopic electrodynamics meets microscopic electronic structure, and it also provides the starting point for introducing the dielectric function.

\subsection{Polarization, electric displacement, and constitutive relations}

\subsection{Linear response in the frequency domain}

\subsection{Anisotropy and the dielectric tensor}

\newpage
\section{From electronic response to the dielectric function}

\subsection{From polarization response to electric susceptibility}

\subsection{The complex dielectric function}

\subsection{Microscopic origin of the imaginary part}

\subsection{Dielectric response in anisotropic materials}

\newpage
\section{Causality and the Kramers--Kronig relations}

\subsection{Causality and analytic structure of response functions}

\subsection{Kramers--Kronig relations for the dielectric function}
\subsection{Static limit and physical interpretation}


\newpage
\section{From the dielectric function to optical properties}

\subsection{Complex refractive index}

\subsection{Absorption coefficient}

\subsection{Reflectivity}

\subsection{Energy-loss function}
